\documentclass[journal]{IEEEtran}
\usepackage{graphicx}
\usepackage[caption=false,font=footnotesize]{subfig}
\usepackage{acronym}
\usepackage{amsmath}
\usepackage{algorithm}
\usepackage{algpseudocode}
\usepackage{multicol}
\usepackage{booktabs}
\usepackage{array}

\usepackage{comment}
\hbadness=10000
\vbadness=10000
% Figures: store all images in figs/ and include them as
% \includegraphics[width=0.45\textwidth]{\paperplotdir/<name>.png}
\graphicspath{{figs/}}
\newcommand{\paperplotdir}{figs}
\newcommand{\panelwidth}{0.32\textwidth} % three panels across a two-column page
\newcommand{\quarterwidth}{0.245\textwidth} % four panels across (Fig. 1)
\usepackage{xcolor}
\newcommand{\edgeshareplot}{%
\begingroup\setlength{\unitlength}{\quarterwidth}%
\begin{picture}(1,0.680000)
\put(0,0){\includegraphics[width=\quarterwidth]{\paperplotdir/Reliability - HRLLC Edge Placement per Round.png}}%
\put(0,0.639){\color{white}\rule{\quarterwidth}{0.042\unitlength}}%
\put(0.5,0.661){\makebox(0,0){\fontsize{4.1}{4.8}\selectfont HRLLC local/edge execution per round}}%
\put(0,0.100){\color{white}\rule{0.037\unitlength}{0.490\unitlength}}%
\put(0.017,0.350){\makebox(0,0){\rotatebox{90}{\fontsize{4.2}{5}\selectfont Local/edge execution share (\%)}}}%
\end{picture}\endgroup}


\usepackage{amssymb}   % or \usepackage{amsfonts}

\usepackage[hidelinks]{hyperref}
\usepackage{aliascnt}

\renewcommand*{\appendixname}{Appendix}
\renewcommand*{\appendixautorefname}{Appendix}

\providecommand*{\algorithmautorefname}{Algorithm}
\providecommand*{\figureautorefname}{Figure}
\providecommand*{\subfigureautorefname}{Figure}
\providecommand*{\Autoref}[1]{%
  \begingroup
  \def\sectionautorefname{Section}%
  \def\subsectionautorefname{Subsection}%
  \def\subsubsectionautorefname{Subsubsection}%
  \def\algorithmautorefname{Algorithm}%
  \def\figureautorefname{Figure}%
  \def\subfigureautorefname{Figure}%
 ~\autoref{#1}%
  \endgroup
}
\newcommand{\Subfigref}[2]{\mbox{\autoref{#1}\subref{#2}}}
\newcommand{\Autoref}[1]{Eq.~\Autoref{#1}}
\newcommand{\Eqsref}[2]{Eqs.~\Autoref{#1}--\Autoref{#2}}
\newcommand{\Eqtworef}[2]{Eqs.~\Autoref{#1} and~\Autoref{#2}}

\renewcommand{\bottomfraction}{0.9}
\renewcommand{\textfraction}{0.05}
\renewcommand{\floatpagefraction}{0.8}
%-------------------------- Acronyms ------------------------
\newacro{iot}[IoT]{Internet of Things}
\newacro{bs}[BS]{base station}
\newacro{uav}[UAV]{unmanned aerial vehicle}
\newacro{mec}[MEC]{Mobile Edge Computing}
\newacro{sinr}[SINR]{signal-to-interference-plus-noise ratio}
\newacro{rl}[RL]{Reinforcement Learning}
\newacro{drl}[DRL]{Deep Reinforcement Learning}
\newacro{ppo}[PPO]{Proximal Policy Optimization}
\newacro{a2c}[A2C]{Advantage Actor-Critic}
\newacro{dqn}[DQN]{Deep Q-Network}
\newacro{ddqn}[DDQN]{Double Deep Q-Network}
\newacro{cpu}[CPU]{Central Processing Unit}
\newacro{urllc}[URLLC]{Ultra-Reliable Low-Latency Communications}
\newacro{mmtc}[mMTC]{massive Machine-Type Communications}
\newacro{embb}[eMBB]{enhanced Mobile Broadband}
\newacro{hrllc}[HRLLC]{hyper-reliable and low-latency communication}
\newacro{mc}[MC]{massive communication}
\newacro{ic}[IC]{immersive communication}
\newacro{e2e}[E2E]{end-to-end}
\newacro{sd}[SD]{standard deviation}
\newacro{cdf}[CDF]{cumulative distribution function}

\hyphenation{op-tical net-works semi-conduc-tor}

\begin{document}

%\title{Deep Reinforcement Learning for Hyper-Reliable and Low-Latency Communication in 6G Massive IoT Edge--Cloud Environments}

%\title{Deep Reinforcement Learning for Hyper-Reliable and Low-Latency Communication in 6G Edge--Cloud Environments}

%\title{Service-Differentiated Deep Reinforcement Learning Offloading for Hyper-Reliable Low-Latency Communication in the 6G Computing Continuum}

%\title{Deep Reinforcement Learning for Service-Differentiated HRLLC Offloading Across the 6G Edge-Cloud Continuum}

%\title{Reliability-Aware Deep Reinforcement Learning for Hyper-Reliable Low-Latency Offloading in the 6G Computing Continuum}


\title{Hyper-Reliable Low-Latency Task Offloading in the 6G Computing Continuum Using Deep Reinforcement Learning}

\author{Afzal Badshah
\thanks{Afzal Badshah is with the Department of Software Engineering, University of Sargodha, Sargodha, Pakistan (e-mail: afzalbadshahphd@gmail.com; afzal.badshah@uos.edu.pk).}}%

\maketitle

\begin{abstract}
6G networks needs to support heterogeneous vehicular and urban sensing workloads under changing connectivity and limited onboard resources. This makes timely task execution difficult under massive connectivity.
%
However, an execution tier suitable for routine sensing may not complete urgent task within deadline. These workloads require reliable completion of real-time tasks, battery preservation for routine sensing devices, and scalable computing capacity for data-intensive processing under increasing demand.
%
To address these challenges, we propose a reliability-aware offloading framework for the 6G computing continuum. A \ac{ddqn} selects feasible execution tiers across onboard, ground-edge, aerial, and cloud resources. The service-differentiated rewards prioritise latency, device battery preservation, and computing adequacy while penalising missed application deadlines.
%
We evaluated task offloading in our custom-built \textit{6G-UrbanNetSim} by comparing policies and delay ablations. In the main simulation campaign, \ac{ddqn} met the 100\% application deadline in all 17,105 \ac{hrllc} attempts, with a pooled 99th-percentile completion time of 1.062~ms and a mean task energy efficiency of 2.900~Mbit/J. In a separate matched ablation, adding the \acs{hrllc} latency term reduced mean per-seed completed-task P99 by 54.1\%, while mean task energy efficiency fell from 3.846 to 3.340~Mbit/J. 
\end{abstract}

\begin{IEEEkeywords}
6G, edge computing, deep reinforcement learning, hyper-reliable low-latency communication, task offloading
\end{IEEEkeywords}

\acresetall

\IEEEpeerreviewmaketitle


\section{Introduction}
\label{sec:introduction}

Connected vehicles and urban sensing depend on timely computing to turn collected data into useful actions. The IMT-2030 framework provides \ac{hrllc} as a 6G usage scenario and gives air-interface research targets of 0.1--1~ms latency and success probabilities ranging from \(1-10^{-5}\) to \(1-10^{-7}\)~\cite{ITU_M2160}.
%
These objectives motivate task offloading that considers where data are processed as well as how they are transmitted~\cite{Liu2017LatencyReliability}.
Across the computing continuum, task completion also depends on access interruptions, queueing, and processing at the selected computing node.
Reliable application operation therefore requires task offloading that connects communication conditions with the available computing resources~\cite{ETSI_MEC002_2025}.

The computing continuum provides execution tiers (onboard devices, ground-edge servers, aerial resources, and  cloud) for meeting these requirements~\cite{Dai2024UAVAssisted}.
Processing tasks at nearby edge servers reduce application latency, though the benefit depends on how wireless connectivity changes as vehicles move.~\cite{ETSI_MEC002_2025}.
Selecting computation tier decision depends on balancing communication delay, available computing resources, and device energy as vehicles move and workload increases~\cite{RW05,RW08,Batewela2020RiskSensitive}.
\Ac{drl} offers an adaptive approach to this task-offloading decision by learning from network conditions and execution outcomes. Its reward reflect the service requirements that determine whether the offloading is useful, motivating service-differentiated rewards across the computing continuum.

The IMT-2030 framework differentiates \ac{hrllc}, \ac{mc}, and \ac{ic}, reflecting different latency, reliability, connectivity, and data-delivery needs~\cite{ITU_M2160}.
In this study, these classes represent alert and control tasks, sensing and telemetry, and data-intensive workloads, respectively.
An offloading decision suitable for routine telemetry may therefore be unsuitable for a real-time task, while moving processing closer to a device can change energy consumption and execution cost. %Evaluating service-differentiated rewards consequently requires reproducible multi-tier benchmarking, service-differentiated task offloading, and joint latency--resource assessment, with \ac{hrllc} tail latency and deadline success examined alongside the resources used.

Existing research provides important foundations for addressing these needs. Learning-based communication has addressed vehicular link reliability and transmission latency, while task-offloading research has incorporated reliability, latency, and energy into reward design~\cite{RW01,RW08}.
Service-aware \ac{drl} has also supported differentiated resource management across radio-access and edge domains~\cite{RW09}.
The challenge is to establish, through reproducible multi-tier benchmarking, how service-differentiated task offloading affects \ac{hrllc} tail latency and application deadline success across the computing continuum as device load increases.
Joint latency--resource assessment must connect these outcomes with execution-tier choices, cost, and energy efficiency, while evaluating the contribution of latency reward shaping.

This study addresses these connected needs through reliability-aware \ac{drl} for task offloading in the 6G computing continuum.
The proposed \textit{6G-UrbanNetSim} platform provides reproducible multi-tier benchmarking under a device load that ramps up from 1 to 200. Service-differentiated rewards guide task-offloading decisions by prioritising \ac{hrllc} latency and penalising deadline misses. Joint latency--resource assessment evaluates offloading policies and the role of \ac{hrllc} latency reward shaping through a matched ablation. We examine how offloading decisions relate to \ac{hrllc} tail latency and application deadline success, while also comparing execution cost and energy efficiency.

The main contributions of this paper are as follows:
\begin{enumerate}
  \item We develop \textit{6G-UrbanNetSim}, a multi-tier edge--cloud simulation platform for vehicular task offloading across heterogeneous execution tiers.
  The platform supports reproducible benchmarking of offloading policies under increasing device loads in the 6G computing continuum.
  \item We formulate a reliability-aware offloading framework that incorporates service-differentiated rewards into execution-tier selection.
  The reward design prioritises \ac{hrllc} latency and penalises deadline misses to guide task offloading across the computing continuum.
  \item We conduct a comparative evaluation of four deep \ac{rl} and two baseline offloading policies using \ac{hrllc} tail latency, deadline success, execution-tier selection, cost, and energy efficiency.
  We additionally retrain \ac{ddqn} without latency reward shaping to examine changes in tail latency, offloading decisions, and energy efficiency.
\end{enumerate}

The rest of this paper is organised as follows.
\Autoref{sec:related-work} reviews related work.
\Autoref{sec:methodology} presents the methodology.
\Autoref{sec:evaluation} evaluate the proposed model and discusses the results.
Finally, \Autoref{sec:conclusion} concludes the paper.








\begin{figure*}[!t]
  \centering
  \includegraphics[width=\textwidth]{figs/hrllc_architecture.png}
  \caption{Architecture of the proposed reliability-aware offloading framework in \textit{6G-UrbanNetSim}}
  \label{fig:hrllc-architecture}
\end{figure*}


\section{Related Work}
\label{sec:related-work}

Existing studies relevant to this work address three main areas: reliability-oriented optimisation at the physical and link layers, \ac{drl}-based task offloading in \ac{mec}, and service-differentiated resource management in 6G networks.

%\subsection{Reliability-Oriented Physical and Link-Layer Optimisation}

\ac{urllc} performance has been pursued primarily through physical-layer and radio-access techniques. \ac{drl}-based scheduling optimised vehicular optical camera communication links to satisfy latency and reliability targets under time-varying channel fading~\cite{RW01}, while scheduling strategies within Open RAN architectures explored trade-offs between throughput and \ac{urllc} compliance under heterogeneous traffic~\cite{RW02}. These single-link formulations demonstrated that learning-based agents can converge to policies satisfying strict delay bounds, but the optimisation scope remains confined to the radio access interface. Extending the scope to spatial and spectral resource control, joint trajectory planning and NOMA-aided beamforming for \ac{uav}-assisted integrated sensing, communication, and computation networks incorporated \ac{urllc} constraints into multi-dimensional flight-path optimisation~\cite{RW03}. Reconfigurable intelligent surfaces were further combined with cell-free massive MIMO to tighten finite-blocklength error bounds through joint beamforming and blocklength control~\cite{RW04}. By coupling spatial, spectral, and coding degrees of freedom, these physical-layer designs approach the reliability floor for individual communication links. The challenge is that these formulations do not see computation-offloading decisions across a multi-tier edge--cloud continuum, where \ac{e2e} latency additionally depends on queuing, compute-processing, and backhaul delay components that lie beyond the radio link.

%\subsection{DRL-Based Task Offloading}

\ac{drl} has been applied to task offloading to manage the combinatorial complexity of multi-user, multi-server scheduling in edge computing environments. Mean-field modelling reduced the state-space dimensionality of dense vehicular offloading by representing aggregate device behaviour rather than tracking individual interactions~\cite{RW05}, while resource allocation in 6G edge-intelligence architectures adopted \ac{drl} to coordinate heterogeneous compute nodes for joint resource allocation and task offloading under dynamic workload conditions~\cite{RW06,Wu2024CloudEdgeEnd}. A reliability-oriented formulation introduced context-aware adaptive reward weighting to balance multiple offloading objectives, including task-completion probability and latency, by adjusting the reward signal according to instantaneous network state~\cite{RW08,Hao2025ReliabilityUAV}. These studies advanced offloading intelligence beyond static or heuristic offloading strategies and reduced mean task-completion latency under high device density. Prior studies have used task priority and mobility to guide vehicular offloading~\cite{Badshah2026PMAQL} and have incorporated priority-dependent latency weights into learned policies~\cite{Uddin2024Prioritized,Barai2026PriorityIoT}. We examine how adding an \acs{hrllc} latency term changes execution-tier selection, completion-time tails, and task energy efficiency when admission rules, spillover, and the deadline-miss penalty are held fixed. We report completed-task P99 to characterise the latency tail and the fraction of all attempts completed within the 100~ms application deadline to measure deadline success.

%\subsection{Resource Management in 6G Networks}

Emerging 6G architectures require resource management that differ among heterogeneous service classes defined under the IMT-2030 framework. Quantum \ac{drl} was applied to satellite--air--ground integrated networks to support \ac{urllc} alongside other traffic types, using a digital-twin mirror of network state to reduce exploration overhead during agent convergence~\cite{RW07}. Transfer-learning-augmented \ac{drl} managed intra- and inter-slice resources across RAN and edge domains to satisfy per-slice quality targets under varying load conditions without retraining from scratch~\cite{RW09}. SLA-driven xApps within O-RAN architectures employed \ac{drl} to enforce per-flow QoS requirements and maintain service-level agreements under dynamic traffic patterns~\cite{RW10}. These approaches address service differentiation at the network-slicing or radio-resource allocation layer and demonstrate that learning-based methods can enforce class-specific guarantees at the scheduling level. They do not, however, extend to computation offloading across edge, aerial, and cloud tiers, nor do they model the interaction between offloading decisions and sub-millisecond tail-latency enforcement for hyper-reliable traffic in the computing continuum.

%\subsection{Research Gap}

Physical-layer reliability, execution-tier selection, and service differentiation affect different parts of task completion. Their interaction motivates measuring whether a reliability-aware offloading reward changes \ac{hrllc} tail latency, application deadline success, and resource use across heterogeneous execution tiers. The present work evaluates these outcomes under increasing device loads and uses a matched \ac{ddqn} reward ablation to examine the latency--efficiency trade-off.



\section{Methodology}
\label{sec:methodology}

The proposed framework coordinates task offloading across onboard, ground-edge, aerial, and cloud computing resources, as shown in \Autoref{fig:hrllc-architecture}. 

\subsection{System and Task Model}
\label{subsec:system-task-model}

Each device (\(d(i)\)) generates task (\(i\)) with payload (\(L_i\)), computational work (\(W_i\)), memory demand (\(M_i\)), and storage demand (\(V_i\)). The service class (\(c_i\)) determines deadline (\(\delta_i\)), so the task is represented by \(\boldsymbol{\tau}_i=(L_i,W_i,M_i,V_i,\delta_i,c_i)\). \acs{hrllc}, \acs{mc}, and \acs{ic} represent alert and control, sensing and telemetry, and data-intensive workloads, respectively.

The policy action set (\(\mathcal A\)) contains four computing tiers $\mathcal A=\{\mathrm{LOCAL},\!\mathrm{MEC},\!\mathrm{UAV},\!\mathrm{CLOUD}\}$. The requested action (\(a_i\)) may differ from the executed action (\(\tilde a_i\)) because of runtime spillover.

Computing node (\(j\)) provides capacity (\(C_{jr}\)) for resource (\(r\in\{\mathrm{cpu},\mathrm{mem},\mathrm{sto}\}\)). Reservations (\(d_{ir}\)) use cycles per second for CPU and bytes for memory and storage. The requested compute rate (\(\rho_i\)), defined by \(\rho_i=d_{i,\mathrm{cpu}}\), therefore differs from computational work \(W_i\).

Each available action (\(a\)) maps to a computing node (\(j_i(a)\)), with \(j_i(\mathrm{LOCAL})=d(i)\). Its task-associated capacities (\(C_{iar}\)) and committed resources (\(U_{iar}\)) satisfy \(C_{iar}=C_{j_i(a),r}\) and \(U_{iar}=U_{j_i(a),r}\). The clipping operator (\(\operatorname{clip}(x,l,h)\)) restricts \(x\) to \([l,h]\). For each remote resource with positive capacity, the projected fraction is \(v_{iar}=\operatorname{clip}((U_{iar}+d_{ir})/C_{iar},0,1)\). The remote resource set (\(\mathcal R_{\mathrm{rem}}\)) is \(\mathcal R_{\mathrm{rem}}=\{\mathrm{cpu},\mathrm{mem},\mathrm{sto}\}\). Current local \acs{cpu} utilization (\(u_{ia}^{\mathrm{cpu}}\)) and remote resource fractions determine projected action utilization (\(\widehat u_{ia}\)):

\begin{equation}
\label{eq:projected-utilization}
\widehat u_{ia}=
\begin{cases}
u_{ia}^{\mathrm{cpu}}
+\dfrac{\rho_i}{C_{ia,\mathrm{cpu}}},
&\begin{gathered}a=\mathrm{LOCAL},\\
C_{ia,\mathrm{cpu}}>0\end{gathered},\\[2mm]
\displaystyle\max_{r\in\mathcal R_{\mathrm{rem}}}v_{iar},
&\begin{gathered}a\neq\mathrm{LOCAL},\\
C_{iar}>0\ \forall r\in\mathcal R_{\mathrm{rem}}\end{gathered},\\
1,&\text{otherwise},
\end{cases}
\end{equation}

The projected load is used in the policy observation, spillover rule, and overload calculation.

Local eligibility (\(\chi_i^{\mathrm{loc}}\)) requires enabled computing and sufficient battery energy when a finite battery is specified. The local computing indicator (\(I_i^{\mathrm{cpu}}\)) and memory indicator (\(I_i^{\mathrm{mem}}\)) check whether the requested reservations fit the remaining computing and memory capacities. For a remote action, the candidate set (\(\mathcal J_{ia}^{\mathrm{adm}}\)) contains operational, reachable servers satisfying \(d_{ir}\leq C_{jr}-U_{jr}\) for every required resource. The indicator (\(\mathbf1\{\cdot\}\)) equals one when its condition is true. These conditions determine the physical mask (\(m_{ia}\)):

\begin{equation}
\label{eq:physical-feasibility}
m_{ia}=
\begin{cases}
\chi_i^{\mathrm{loc}}
I_i^{\mathrm{cpu}}I_i^{\mathrm{mem}},
&a=\mathrm{LOCAL},\\
\mathbf1\{\mathcal J_{ia}^{\mathrm{adm}}\neq\varnothing\},
&a\neq\mathrm{LOCAL}.
\end{cases}
\end{equation}

The physical mask in~\Autoref{eq:physical-feasibility} defines the feasible action set (\(\mathcal A_i^{\mathrm f}\)) through \(\mathcal A_i^{\mathrm f}=\{a\in\mathcal A:m_{ia}=1\}\). 


\subsection{Communication and Task Completion}
\label{subsec:communication-completion}

An admitted remote action follows route (\(\mathcal P_i\)), whose hops are indexed by (\(\ell\)). Hop bandwidth (\(B_\ell\)), efficiency (\(\eta_\ell\)), spatial gain (\(\chi_\ell\)), and SINR (\(\Gamma_\ell\)) determine the achievable rate. The outage indicator (\(o_\ell\)) identifies a hop whose received power or \acs{sinr} falls below its decoding threshold. The wireless rate (\(R_\ell\)) is:

\begin{equation}
\label{eq:wireless-rate}
R_\ell
=
(1-o_\ell)\eta_\ell B_\ell\chi_\ell
\log_2(1+\Gamma_\ell).
\end{equation}

The configured spatial-stream count (\(n_\ell^{\mathrm{str}}\)) and multiplexing efficiency (\(\eta_{\mathrm{MIMO}}\)) determine the spatial gain through \(\chi_\ell=1+(n_\ell^{\mathrm{str}}-1)\eta_{\mathrm{MIMO}}\). Wired hops use their configured link rates.

Handover, switching, and coverage intervals may overlap before transmission. Access waiting is the maximum of these intervals:

\begin{equation}
\label{eq:access-waiting}
T_i^{\mathrm{acc}}
=
\max\left\{
T_i^{\mathrm{ho}},
T_i^{\mathrm{int}},
T_i^{\mathrm{cov}}
\right\}.
\end{equation}

A task released at time (\(t_i^0\)) has a network start time (\(t_i^{\mathrm{net,s}}\)) given by \(t_i^{\mathrm{net,s}}=t_i^0+T_i^{\mathrm{acc}}\). This construction counts overlapping interruptions once.

For packet index (\(p\)), hop \(\ell\), and transmission-attempt index (\(b\)), the on-wire size (\(L_{ip\ell b}^{\mathrm w}\)) includes packet headers and link overhead. For an available hop with rate from~\Autoref{eq:wireless-rate}, the serialization duration (\(T_{ip\ell b}^{\mathrm{tx}}\)) is \(T_{ip\ell b}^{\mathrm{tx}}=8L_{ip\ell b}^{\mathrm w}/R_\ell\). Hop distance (\(D_\ell\)) and propagation speed (\(v_\ell\)) give the propagation duration (\(T_\ell^{\mathrm{prop}}\)) through \(T_\ell^{\mathrm{prop}}=D_\ell/v_\ell\).

Packet readiness and shared-resource availability determine the transmission start time (\(t_{ip\ell b}^{\mathrm s}\)). The resulting arrival time (\(t_{ip\ell b}^{\mathrm{arr}}\)) and updated resource availability follow:

\begin{equation}
\label{eq:shared-link-scheduling}
\begin{aligned}
t_{ip\ell b}^{\mathrm s}
&=\max(t_{ip\ell b}^{\mathrm{rdy}},F_\ell),\\
t_{ip\ell b}^{\mathrm{arr}}
&=t_{ip\ell b}^{\mathrm s}
+T_{ip\ell b}^{\mathrm{tx}}
+T_\ell^{\mathrm{prop}},\\
F_\ell
&\leftarrow
t_{ip\ell b}^{\mathrm s}+T_{ip\ell b}^{\mathrm{tx}}.
\end{aligned}
\end{equation}

The shared-link schedule in~\Autoref{eq:shared-link-scheduling} gives the queueing interval \(t_{ip\ell b}^{\mathrm s}-t_{ip\ell b}^{\mathrm{rdy}}\). Successful reception makes the packet ready for its next hop. 

Packet termination time (\(t_{ip}^{\mathrm{term}}\)) records successful delivery or failure. The final packet termination determines the task’s elapsed network duration (\(T_i^{\mathrm{net}}\)):

\begin{equation}
\label{eq:network-duration}
T_i^{\mathrm{net}}
=
\max_p t_{ip}^{\mathrm{term}}
-t_i^{\mathrm{net,s}}.
\end{equation}

The timeline already includes queueing, retries, and outages; these intervals are not added again.

The selected computing node processes the task at its full compute rate (\(f_{i,\tilde a_i}\)), with \(f_{i,\tilde a_i}=C_{i,\tilde a_i,\mathrm{cpu}}\). Existing \acs{cpu} utilization (\(u_i^{-}\)) increases processing duration through bounded congestion. The minimum duration (\(T_{\min}\)) is in seconds, with \(0<h_{\min}\leq1\). With computation indicator (\(d_i^{\mathrm{exec}}\)), processing duration (\(T_i^{\mathrm{proc}}\)) is:

\begin{equation}
\label{eq:processing-duration}
T_i^{\mathrm{proc}}
=
d_i^{\mathrm{exec}}
\frac{
\max\!\left(T_{\min},W_i/f_{i,\tilde a_i}\right)
}{
\max(1-u_i^{-},h_{\min})
}.
\end{equation}

Local execution has \(d_i^{\mathrm{exec}}=1\). Remote computation occurs only after full-payload delivery.

Complete task latency is the sum of access waiting, network duration, and processing duration:

\begin{equation}
\label{eq:complete-task-latency}
T_i
=
T_i^{\mathrm{acc}}
+
T_i^{\mathrm{net}}
+
T_i^{\mathrm{proc}}.
\end{equation}

Local execution has zero access and network durations. The timely-completion indicator (\(z_i\)) is \(z_i=\mathbf1\{\text{application completed}\}\mathbf1\{T_i\leq\delta_i\}\), where application completion includes the final resource-admission check. The complete latency in~\Autoref{eq:complete-task-latency} also determines the waiting interval used in device-energy accounting.

\subsection{Energy and Cost Accounting}
\label{subsec:energy-cost-accounting}

Computing node power consumption depends on \acs{cpu} utilization (\(u\)). Idle power (\(P_j^{\mathrm{idle}}\)), dynamic-power coefficient (\(P_j^{\mathrm{dyn}}\)), maximum power (\(P_j^{\max}\)), and utilization exponent (\(\xi_j\)) define the configured power curve (\(P_j(u)\)) through \(P_j(u)=\min(P_j^{\max},P_j^{\mathrm{idle}}+P_j^{\mathrm{dyn}}u^{\xi_j})\).

For local execution, device active power (\(P_{d(i)}^{\mathrm{act}}\)) determines computing energy. For remote execution, \acs{cpu} utilization after admission (\(u_i^{+}\)) is compared with its preceding value \(u_i^{-}\). The positive-part operator (\([x]_+\)) is defined by \([x]_+=\max(x,0)\). The task-attributed computing energy (\(E_i^{\mathrm{IT}}\)) is:

\begin{equation}
\label{eq:computing-energy}
E_i^{\mathrm{IT}}
=
\begin{cases}
P_{d(i)}^{\mathrm{act}}T_i^{\mathrm{proc}},
&\tilde a_i=\mathrm{LOCAL},\\[1mm]
\left[
P_j(u_i^{+})-P_j(u_i^{-})
\right]_+
T_i^{\mathrm{proc}},
&\tilde a_i\neq\mathrm{LOCAL},
\end{cases}
\end{equation}

The computing node has \(\mathrm{PUE}_j\geq1\). Remote facility overhead is therefore \(E_i^{\mathrm{fac}}=(\mathrm{PUE}_j-1)E_i^{\mathrm{IT}}\). The configured local device profiles have zero facility overhead. Incremental remote accounting avoids charging the shared computing node’s idle baseline to every task.

Device transmission energy (\(E_i^{\mathrm{tx}}\)) uses electrical transmit power (\(P_{d(i)}^{\mathrm{tx}}\)) and reserved access airtime (\(T_i^{\mathrm{air}}\)). For an available device power profile, \(E_i^{\mathrm{tx}}=P_{d(i)}^{\mathrm{tx}}T_i^{\mathrm{air}}\); otherwise, the packet transmission-energy record supplies the value. Originating-device reception energy (\(E_i^{\mathrm{rx}}\)) is zero for the modeled upload.

Device idle power (\(P_{d(i)}^{\mathrm{idle}}\)) determines waiting energy (\(E_i^{\mathrm{wait}}\)) through \(E_i^{\mathrm{wait}}=P_{d(i)}^{\mathrm{idle}}[T_i-T_i^{\mathrm{air}}]_+\). This waiting interval includes access waiting and the time spent awaiting transmission or computation outcomes. Failed delivery can therefore consume communication and waiting energy even though computing energy is zero.

Infrastructure network energy (\(E_i^{\mathrm{net,inf}}\)) includes reception on the first access hop and transmit, receive, and transport energy on subsequent hops. Originating-device transmission is accounted for separately. The task-energy ledger is:

\[
\mathbf e_i=
\left(
E_i^{\mathrm{tx}},
E_i^{\mathrm{rx}},
E_i^{\mathrm{wait}},
E_i^{\mathrm{net,inf}},
E_i^{\mathrm{IT}},
E_i^{\mathrm{fac}}
\right)^{\mathsf T}.
\]

The local-execution indicator (\(\lambda_i^{\mathrm{loc}}\)), defined by \(\lambda_i^{\mathrm{loc}}=\mathbf1\{\tilde a_i=\mathrm{LOCAL}\}\), determines which computing components draw energy from the originating device. Its ownership vector (\(\boldsymbol\omega_i\)) is \(\boldsymbol\omega_i=(1,1,1,0,\lambda_i^{\mathrm{loc}},\lambda_i^{\mathrm{loc}})^{\mathsf T}\). Device battery expenditure (\(E_i^{\mathrm{dev}}\)) is:

\begin{equation}
\label{eq:device-battery-energy}
E_i^{\mathrm{dev}}
=
\boldsymbol\omega_i^{\mathsf T}\mathbf e_i.
\end{equation}

For a finite device battery, pre-task energy (\(B_i^{\mathrm{bef}}\)) is reduced to the post-task energy (\(B_i^{\mathrm{aft}}\)) through \(B_i^{\mathrm{aft}}=[B_i^{\mathrm{bef}}-E_i^{\mathrm{dev}}]_+\). The expenditure in~\Autoref{eq:device-battery-energy} determines this debit and the battery state exposed to subsequent decisions.

The all-ones vector (\(\mathbf1\)) includes every ledger component in total task-attributed energy (\(E_i^{\mathrm{tot}}\)):

\begin{equation}
\label{eq:total-task-energy}
E_i^{\mathrm{tot}}
=
\mathbf1^{\mathsf T}\mathbf e_i.
\end{equation}

The total in~\Autoref{eq:total-task-energy} attributes task communication, waiting, and computing energy. \acs{uav} hovering, propulsion, and shared idle baselines are tracked separately.

For component index (\(q\)), operator cost (\(C_i^{\mathrm{op}}\)) pairs each measured billing quantity (\(\nu_{iq}\)) with its configured unit price (\(p_{q,\tilde a_i}\)). The component set (\(\mathcal Q\)) is \(\mathcal Q=\{\mathrm{compute},\mathrm{transport},\mathrm{energy}\}\):

\begin{equation}
\label{eq:operator-cost}
C_i^{\mathrm{op}}
=
\sum_{q\in\mathcal Q}
p_{q,\tilde a_i}\nu_{iq}.
\end{equation}

For remote execution, the reference-vCPU rate (\(F_{\mathrm{ref}}\)) is measured in cycles per second. The compute billing quantity (\(\nu_{i,\mathrm{compute}}\)) is \(\nu_{i,\mathrm{compute}}=d_i^{\mathrm{exec}}W_i/F_{\mathrm{ref}}\), measured in reference-vCPU seconds. Transport quantity (\(\nu_{i,\mathrm{transport}}\)) is the total observed on-wire bytes across the executed route. Electricity quantity (\(\nu_{i,\mathrm{energy}}\)) is \(\nu_{i,\mathrm{energy}}=E_i^{\mathrm{net,inf}}+E_i^{\mathrm{IT}}+E_i^{\mathrm{fac}}\), measured in joules. 

\subsection{Reliability-Aware Task Offloading}
\label{subsec:reliability-aware-offloading}

The physical mask excludes actions exceeding resource limits. The operating threshold (\(u_0\)) controls spillover under high projected utilization. Candidates are considered in the configured order of Local, \acs{mec}, \acs{uav}, and Cloud, starting at the requested action. Unavailable actions are skipped.

The ordered list (\(\mathcal F_i(a_i)\)) contains physically feasible candidates. Its ordered subset (\(\mathcal G_i(a_i)\)) retains candidates satisfying \(\widehat u_{ia}<u_0\), together with feasible Cloud execution. The executed action is selected as follows:

\begin{equation}
\label{eq:execution-resolution}
\tilde a_i=
\begin{cases}
\operatorname{first}(\mathcal G_i(a_i)),
&\mathcal G_i(a_i)\neq\varnothing,\\
\operatorname{last}(\mathcal F_i(a_i)),
&\mathcal G_i(a_i)=\varnothing,\
\mathcal F_i(a_i)\neq\varnothing,\\
\varnothing,
&\text{otherwise}.
\end{cases}
\end{equation}

Under the resolution rule in~\Autoref{eq:execution-resolution}, the last candidate in the prescribed feasible order executes when every candidate exceeds the operating threshold. Exceeding the soft threshold alone does not cause admission failure. An empty physical action set produces task failure.

Complete latency determines the latency-preference score (\(s_i^{T}\)). The class-dependent preference scale (\(\tau_{c_i}\)), satisfying \(\tau_{c_i}>0\), controls its sensitivity independently of application deadline \(\delta_i\):

\begin{equation}
\label{eq:latency-preference}
s_i^{T}
=
\operatorname{clip}
\left(
\frac{2}{1+T_i/\tau_{c_i}}-1,
-1,1
\right).
\end{equation}

For action \(a\), current resource load (\(u_{ia}\)) determines free headroom (\(h_{ia}\)) through \(h_{ia}=\operatorname{clip}(1-u_{ia},0,1)\). Current load is local \acs{cpu} utilization for onboard execution and peak \acs{cpu}, memory, or storage utilization for remote execution. Compute adequacy (\(\kappa_{ia}\)), defined by \(\kappa_{ia}=\operatorname{clip}(C_{ia,\mathrm{cpu}}/\rho_i,0,1)\), compares full compute capacity with positive requested compute rate. Absent capacity gives zero adequacy. These observations determine the compute score (\(s_i^{K}\)):

\begin{equation}
\label{eq:compute-adequacy}
s_i^{K}
=
\frac{
h_{i,\tilde a_i}
+
\kappa_{i,\tilde a_i}
}{2}.
\end{equation}

The score combines current headroom and task-relative capacity. Projected saturation of the requested tier enters separately through overload penalty (\(p_i^{O}\)):

\begin{equation}
\label{eq:overload-penalty}
p_i^{O}
=
\left[
\operatorname{clip}
\left(
\frac{\widehat u_{i,a_i}-u_0}{1-u_0},
0,1
\right)
\right]^2.
\end{equation}

%This penalty remains associated with the requested action even when spillover changes the executed tier.

The device battery capacity (\(B_{d(i)}^{\max}\)) is measured in joules. Its pre-debit battery ratio (\(\beta_i\)) is \(\beta_i=\operatorname{clip}(B_i^{\mathrm{bef}}/B_{d(i)}^{\max},0,1)\) when \(B_{d(i)}^{\max}>0\), and \(\beta_i=1\) otherwise. The device battery score (\(s_i^{B}\)) combines this ratio with realized device expenditure. A positive class-dependent energy normalizer (\(N_{c_i}^{E}\)), measured in joules, gives the normalized device expenditure (\(p_i^{E}\)) through \(p_i^{E}=\operatorname{clip}(E_i^{\mathrm{dev}}/N_{c_i}^{E},0,3)\). The score is \(s_i^{B}=\operatorname{clip}(\beta_i-p_i^{E}/3,-1,1)\). The executed server’s projected remaining battery ratio (\(\widehat\beta_i^{\mathrm{srv}}\)) is supplied to the reward evaluator. When no projected ratio is supplied, its value is one. The server battery score (\(s_i^{S}\)) is \(s_i^{S}=\operatorname{clip}(\widehat\beta_i^{\mathrm{srv}},-1,1)\).

The cost penalty (\(p_i^{C}\)) uses a positive monetary normalizer (\(N_{c_i}^{C}\)) through \(p_i^{C}=\operatorname{clip}(C_i^{\mathrm{op}}/N_{c_i}^{C},0,3)\). The remaining coverage interval (\(T_i^{\mathrm{dwell}}\)) is measured in seconds. The mobility-sensitive action set (\(\mathcal A_{\mathrm{mob}}\)) contains \acs{mec} and \acs{uav} execution. The mobility penalty (\(p_i^{M}\)) is \(p_i^{M}=\mathbf1\{\tilde a_i\in\mathcal A_{\mathrm{mob}},\ T_i^{\mathrm{dwell}}>0,\ T_i>T_i^{\mathrm{dwell}}\}\). For the executed tier, define clipped current load (\(\bar u_i^{\mathrm{cur}}\)) through \(\bar u_i^{\mathrm{cur}}=\operatorname{clip}(u_{i,\tilde a_i},0,1)\) and projected load (\(\bar u_i^{\mathrm{proj}}\)) through \(\bar u_i^{\mathrm{proj}}=\operatorname{clip}(\widehat u_{i,\tilde a_i},0,1)\). The shared-node indicator (\(\sigma_i^{\mathrm{sh}}\)) is \(\sigma_i^{\mathrm{sh}}=\mathbf1\{\tilde a_i\neq\mathrm{LOCAL}\}\). The shared-capacity penalty (\(p_i^{X}\)) is \(p_i^{X}=\sigma_i^{\mathrm{sh}}\operatorname{clip}((\bar u_i^{\mathrm{proj}})^2-(\bar u_i^{\mathrm{cur}})^2,0,1)\).

The score group (\(\mathcal K_s=\{T,B,S,K\}\)) contains latency, device-battery, server-battery, and computing scores. The penalty group (\(\mathcal K_p=\{C,M,X,O\}\)) contains cost, mobility, shared-capacity, and overload penalties. Class weights (\(w_{c_i}^{q}\)) define the aggregate score (\(G_i\)) and penalty (\(H_i\)):

\[
\begin{aligned}
G_i
&=\sum_{q\in\mathcal K_s}w_{c_i}^{q}s_i^{q},\\
H_i
&=\sum_{q\in\mathcal K_p}w_{c_i}^{q}p_i^{q}.
\end{aligned}
\]

The class-neutral overload coefficient (\(w^{O}\)) satisfies \(w_{c_i}^{O}=w^{O}\). The normalization component set (\(\mathcal K_Z\)) is \(\mathcal K_Z=\mathcal K_s\cup(\mathcal K_p\setminus\{O\})\). The normalization (\(Z_{c_i}\)) is \(Z_{c_i}=\sum_{q\in\mathcal K_Z}\lvert w_{c_i}^{q}\rvert>0\), excluding the separate overload coefficient. The normalized service objective (\(\Psi_i\)) is:

\begin{equation}
\label{eq:service-objective}
\Psi_i
=
\operatorname{clip}
\left(
\frac{G_i-H_i}{Z_{c_i}},
-1,1
\right).
\end{equation}

The service objective in~\Autoref{eq:service-objective} combines the latency preference from~\Autoref{eq:latency-preference} and compute score from~\Autoref{eq:compute-adequacy} with battery and penalty contributions. Class weights emphasize \acs{hrllc} latency, \acs{mc} device battery preservation, and \acs{ic} computing adequacy. The overload penalty in~\Autoref{eq:overload-penalty} retains pressure associated with the requested tier.

The reward scale (\(R_{\max}\)), satisfying \(R_{\max}>0\), bounds the realized reward (\(r_i\)). The nonnegative coefficient (\(\lambda_{\mathrm{inv}}\)) specifies the invalid-action penalty, and the binary indicator (\(\iota_i^{\mathrm{inv}}\)), with \(\iota_i^{\mathrm{inv}}\in\{0,1\}\), records the reward evaluator’s invalid-action flag. The invalid-action debit (\(p_i^{\mathrm{inv}}\)) is \(p_i^{\mathrm{inv}}=\lambda_{\mathrm{inv}}\iota_i^{\mathrm{inv}}\). Timely completion determines:

\begin{equation}
\label{eq:realized-reward}
r_i=
\begin{cases}
\operatorname{clip}
\left(
R_{\max}\Psi_i-p_i^{\mathrm{inv}},
-R_{\max},R_{\max}
\right),
&z_i=1,\\
-R_{\max},
&z_i=0.
\end{cases}
\end{equation}

Under the reward rule in~\Autoref{eq:realized-reward}, execution failure, final admission failure, and missed deadlines share the failure reward. The \acs{hrllc} ablation removes only its latency contribution from \(G_i\); the original normalization, remaining contributions, and failure reward are retained.

An offloading policy (\(\pi\)) maps observations to action-selection probabilities. Its formal objective is the policy (\(\pi^\star\)) maximizing expected realized task reward (\(\mathbb E_\pi[r_i]\)):

\begin{equation}
\label{eq:policy-objective}
\pi^\star
\in
\arg\max_\pi
\mathbb E_\pi[r_i].
\end{equation}

The learning procedure below approximates the reward objective in~\Autoref{eq:policy-objective} using observable device and infrastructure conditions.

\subsection{\texorpdfstring{\acs{ddqn}}{DDQN}-Based Policy Learning}
\label{subsec:ddqn-policy-learning}

The observation (\(\mathbf s_i\)) combines task/context features (\(\mathbf x_i\)) with action-specific feature blocks (\(\boldsymbol\phi_{ia}\)). For feature index (\(f\)), the fixed normalizer (\(\mathcal N_f\)) transforms a feature value (\(x_f\)) using its positive scale (\(b_f\)): \(\mathcal N_f(x_f)=\operatorname{clip}(x_f/b_f,0,1)\).

The context vector describes task, device, network, battery, and service-class conditions. One feature block is included for each of the four actions. The resulting observation is:

\begin{equation}
\label{eq:observation-encoding}
\mathbf s_i
=
\left[
\mathbf x_i;
\boldsymbol\phi_{i0};
\ldots;
\boldsymbol\phi_{i,|\mathcal A|-1}
\right].
\end{equation}

The complete feature order, block structure, and fixed normalization scales for~\Autoref{eq:observation-encoding} are documented in~{\def\sectionautorefname{Appendix}\autoref{sec:appendix-state-encoding}}. The observation uses observable conditions; counterfactual realized latency, energy, and cost are not supplied as policy inputs.

Online parameters (\(\boldsymbol\theta\)) contain the layer weights (\(\mathbf W_1,\mathbf W_2,\mathbf W_3\)). Bias vectors (\(\mathbf b_1,\mathbf b_2,\mathbf b_3\)) shift the corresponding affine transformations. ReLU activation (\(\sigma\)) forms hidden representations (\(\mathbf h_i^{(1)}\)) and (\(\mathbf h_i^{(2)}\)), which produce the action-value vector (\(\mathbf Q_{\boldsymbol\theta}(\mathbf s_i)\)):

\begin{equation}
\label{eq:action-value-network}
\begin{aligned}
\mathbf h_i^{(1)}
&=\sigma(\mathbf W_1\mathbf s_i+\mathbf b_1),\\
\mathbf h_i^{(2)}
&=\sigma(\mathbf W_2\mathbf h_i^{(1)}+\mathbf b_2),\\
\mathbf Q_{\boldsymbol\theta}(\mathbf s_i)
&=\mathbf W_3\mathbf h_i^{(2)}+\mathbf b_3.
\end{aligned}
\end{equation}

The scalar action value (\(Q_{\boldsymbol\theta}(\mathbf s_i,a)\)) is the action-\(a\) output of the network in~\Autoref{eq:action-value-network}.

For a nonempty feasible set with feasible-action count (\(|\mathcal A_i^{\mathrm f}|\)), greedy action (\(a_i^{\mathrm g}\)) maximizes the predicted value over \(\mathcal A_i^{\mathrm f}\), with a fixed action-index order resolving ties. Training-round index (\(n\)) determines exploration probability (\(\epsilon_n\)). The feasibility-masked policy (\(\pi_n\)) is:

\begin{equation}
\label{eq:masked-policy}
\pi_n(a\mid\mathbf s_i)
=
\frac{\epsilon_n}{|\mathcal A_i^{\mathrm f}|}
+
(1-\epsilon_n)\mathbf1\{a=a_i^{\mathrm g}\},
\qquad a\in\mathcal A_i^{\mathrm f}.
\end{equation}

The initial exploration probability (\(\epsilon_{\mathrm{init}}\)) decreases to the minimum exploration probability (\(\epsilon_{\min}\)) over the decay horizon (\(N_\epsilon\)), following \(\epsilon_n=\epsilon_{\mathrm{init}}+(\epsilon_{\min}-\epsilon_{\mathrm{init}})\min(1,n/N_\epsilon)\). During evaluation, the policy in~\Autoref{eq:masked-policy} uses \(\epsilon_n=0\). Empty-set failures are recorded without selecting an action.

Replay memory stores eligible attempted decisions using the observation, requested action, realized reward, next observation, and terminal indicator. Consequently, replay credits requested action \(a_i\), while its reward reflects execution on \(\tilde a_i\). Failures identified before any feasible policy action exists do not create a policy transition.

Target parameters (\(\boldsymbol\theta^{-}\)) form a periodically synchronized copy of the online network. For next observation (\(\mathbf s_i'\)), the feasible set (\(\mathcal A_i^{\mathrm f\prime}\)) restricts online next-action selection. Its greedy action (\(a_i^{\prime\mathrm g}\)) is selected by the online network and evaluated by the target network when continuation applies. The configured discount factor (\(\gamma\)) and terminal indicator (\(d_i\)) determine learning target (\(y_i\)):

\begin{equation}
\label{eq:learning-target}
y_i=
\begin{cases}
r_i,
&d_i=1
\ \text{or}\
\mathcal A_i^{\mathrm f\prime}=\varnothing,\\
r_i+\gamma
Q_{\boldsymbol\theta^{-}}
(\mathbf s_i',a_i^{\prime\mathrm g}),
&\text{otherwise}.
\end{cases}
\end{equation}

Every task transition in the evaluated \acs{ddqn} implementation is terminal. The target in~\Autoref{eq:learning-target} is therefore the observed task reward, \(y_i=r_i\), and the network estimates the expected immediate reward of each feasible execution option. We report the \acs{dqn} and \acs{ddqn} runs separately.

The gradient-update index (\(k\)) takes values \(k=1,2,\ldots\), with initial online parameters (\(\boldsymbol\theta_1\)). At update \(k\), a replay minibatch (\(\mathcal B_k\)) of size (\(N_b\)) contains transitions sampled uniformly without replacement. Updates occur only after replay warm-up and at the configured attempted-decision interval. The mean-squared learning loss (\(\mathcal L_k\)) is:

\begin{equation}
\label{eq:replay-loss}
\mathcal L_k(\boldsymbol\theta)
=
\frac{1}{|\mathcal B_k|}
\sum_{i\in\mathcal B_k}
\left[
Q_{\boldsymbol\theta}(\mathbf s_i,a_i)-y_i
\right]^2.
\end{equation}

The positive gradient-norm threshold (\(g_{\max}\)) bounds the global norm of the loss gradient (\(\mathbf g_k\)) of~\Autoref{eq:replay-loss}, evaluated at the current parameters: \(\mathbf g_k=\nabla_{\boldsymbol\theta}\mathcal L_k(\boldsymbol\theta_k)\). Its clipped form (\(\widetilde{\mathbf g}_k\)) is \(\widetilde{\mathbf g}_k=\mathbf g_k/\max(1,\|\mathbf g_k\|_2/g_{\max})\).

Adam moment coefficients (\(\mu_1\)) and (\(\mu_2\)) update the first moment (\(\mathbf m_k\)) and second moment (\(\mathbf v_k\)) from this clipped gradient. Starting from the initial first moment (\(\mathbf m_0\)) and initial second moment (\(\mathbf v_0\)), with \(\mathbf m_0=\mathbf v_0=\mathbf0\), their recursions are \(\mathbf m_k=\mu_1\mathbf m_{k-1}+(1-\mu_1)\widetilde{\mathbf g}_k\) and \(\mathbf v_k=\mu_2\mathbf v_{k-1}+(1-\mu_2)\widetilde{\mathbf g}_k^{\,2}\). The bias-corrected first moment (\(\widehat{\mathbf m}_k\)) and second moment (\(\widehat{\mathbf v}_k\)) are \(\widehat{\mathbf m}_k=\mathbf m_k/(1-\mu_1^k)\) and \(\widehat{\mathbf v}_k=\mathbf v_k/(1-\mu_2^k)\).

Learning rate (\(\alpha\)) and numerical stabilizer (\(\varepsilon_{\mathrm A}\)), satisfying \(\varepsilon_{\mathrm A}>0\), determine the parameter update:

\begin{equation}
\label{eq:parameter-update}
\boldsymbol\theta_{k+1}
=
\boldsymbol\theta_k
-
\alpha
\frac{
\widehat{\mathbf m}_k
}{
\sqrt{\widehat{\mathbf v}_k}
+\varepsilon_{\mathrm A}
}.
\end{equation}

All moment operations and the parameter update in~\Autoref{eq:parameter-update} are componentwise. Target parameters are synchronized after each configured gradient-update interval. In the task-terminal experiments, the target does not contain a future-value contribution.

%Training ends after the configured round budget. During evaluation, online parameters remain fixed and the policy selects greedy feasible actions, while physical admission, spillover, communication, and resource accounting continue to determine realized task outcomes.


\begin{algorithm}[t]
\caption{DDQN-Based Reliability-Aware Task Offloading}
\label{alg:ddqn-offloading}
\begin{algorithmic}[1]
\Require task $\boldsymbol{\tau}_i=(L_i,W_i,M_i,V_i,\delta_i,c_i)$, online parameters $\boldsymbol\theta$, target parameters $\boldsymbol\theta^{-}$, replay memory $\mathcal D$, exploration probability $\epsilon_n$
\Ensure executed tier $\tilde a_i$, realized reward $r_i$

\State $\mathbf s_i\leftarrow[\mathbf x_i;\boldsymbol\phi_{i0};\ldots;\boldsymbol\phi_{i,|\mathcal A|-1}]$ \Comment{Eq.~\Autoref{eq:observation-encoding}}

\ForAll{$a\in\mathcal A$} \Comment{Physical feasibility mask}
    \State Compute $m_{ia}$ \Comment{Eq.~\Autoref{eq:physical-feasibility}}
\EndFor
\State $\mathcal A_i^{\mathrm f}\leftarrow\{a\in\mathcal A:m_{ia}=1\}$
\If{$\mathcal A_i^{\mathrm f}=\varnothing$}
    \State Record failure; \textbf{return}
\EndIf

\State $a_i^{\mathrm g}\leftarrow\arg\max_{a\in\mathcal A_i^{\mathrm f}}Q_{\boldsymbol\theta}(\mathbf s_i,a)$ \Comment{Greedy action}
\State Sample $a_i\sim\pi_n(\cdot\mid\mathbf s_i)$ \Comment{Eq.~\Autoref{eq:masked-policy}}

\State Resolve $\tilde a_i$ via $\mathcal G_i(a_i),\mathcal F_i(a_i)$ \Comment{Eq.~\Autoref{eq:execution-resolution}}

\State Execute task on tier $\tilde a_i$:
\State \quad $T_i\leftarrow T_i^{\mathrm{acc}}+T_i^{\mathrm{net}}+T_i^{\mathrm{proc}}$ \Comment{Eq.~\Autoref{eq:complete-task-latency}}
\State \quad $E_i^{\mathrm{dev}}\leftarrow\boldsymbol\omega_i^{\mathsf T}\mathbf e_i$ \Comment{Eq.~\Autoref{eq:device-battery-energy}}

\State $z_i\leftarrow\mathbf1\{\text{completed}\}\cdot\mathbf1\{T_i\leq\delta_i\}$ \Comment{Timely completion}

\State $\Psi_i\leftarrow\operatorname{clip}\!\left(\frac{G_i-H_i}{Z_{c_i}},-1,1\right)$ \Comment{Eq.~\Autoref{eq:service-objective}}

\If{$z_i=1$} \Comment{Eq.~\Autoref{eq:realized-reward}}
    \State $r_i\leftarrow\operatorname{clip}(R_{\max}\Psi_i-p_i^{\mathrm{inv}},-R_{\max},R_{\max})$
\Else
    \State $r_i\leftarrow -R_{\max}$
\EndIf

\State $\mathcal D\leftarrow\mathcal D\cup\{(\mathbf s_i,a_i,r_i,\mathbf s_i',d_i)\}$ \Comment{Replay store}

\If{$|\mathcal D|\geq N_{\mathrm{warm}}$ \textbf{and} update interval reached}
    \State Sample $\mathcal B_k\subset\mathcal D$, $|\mathcal B_k|=N_b$
    \State $y_i\leftarrow r_i\quad\forall\,i\in\mathcal B_k$ \Comment{Terminal; Eq.~\Autoref{eq:learning-target}}
    \State $\mathcal L_k\leftarrow\frac{1}{|\mathcal B_k|}\sum_{i\in\mathcal B_k}[Q_{\boldsymbol\theta}(\mathbf s_i,a_i)-y_i]^2$ \Comment{Eq.~\Autoref{eq:replay-loss}}
    \State $\mathbf g_k\leftarrow\nabla_{\boldsymbol\theta}\mathcal L_k;\quad\widetilde{\mathbf g}_k\leftarrow\mathbf g_k/\max(1,\|\mathbf g_k\|_2/g_{\max})$
    \State $\boldsymbol\theta_{k+1}\leftarrow\boldsymbol\theta_k-\alpha\,\widehat{\mathbf m}_k/(\sqrt{\widehat{\mathbf v}_k}+\varepsilon_{\mathrm A})$ \Comment{Eq.~\Autoref{eq:parameter-update}}
    \If{sync interval reached}
        \State $\boldsymbol\theta^{-}\leftarrow\boldsymbol\theta$
    \EndIf
\EndIf

\State $\epsilon_n\leftarrow\epsilon_{\mathrm{init}}+(\epsilon_{\min}-\epsilon_{\mathrm{init}})\min(1,n/N_\epsilon)$ \Comment{Decay}
\end{algorithmic}
\end{algorithm}





\subsection{Algorithm}
\autoref{alg:ddqn-offloading} consolidates the decision pipeline. For each arriving task, the agent constructs a normalized observation~\Autoref{eq:observation-encoding}, masks physically infeasible tiers~\Autoref{eq:physical-feasibility}, and selects an action under the $\epsilon$-greedy feasibility-masked policy~\Autoref{eq:masked-policy}. Runtime spillover resolves the requested tier to an executed tier~\Autoref{eq:execution-resolution}. The environment returns complete latency~\Autoref{eq:complete-task-latency} and device energy~\Autoref{eq:device-battery-energy}, from which the class-weighted service objective~\Autoref{eq:service-objective} and realized reward~\Autoref{eq:realized-reward} are computed. The transition is stored in replay memory. At each configured interval, a minibatch gradient step with norm clipping and Adam optimization~\Autoref{eq:parameter-update} updates the online network, and the target network is synchronized periodically.







\begin{table}[t]
\centering
\caption{Experimentation Configration}
\label{tab:experimental-configuration}
\footnotesize
\setlength{\tabcolsep}{3pt}
\renewcommand{\arraystretch}{1.1}
\begin{tabular}{@{}
>{\raggedright\arraybackslash}p{\dimexpr0.46\columnwidth-\tabcolsep\relax}
>{\raggedright\arraybackslash}p{\dimexpr0.54\columnwidth-\tabcolsep\relax}
@{}}
\toprule
Parameter & Value or range \\
\midrule
Area; ground \acs{bs} sites
& $10\times10$~km; 81 \\
\acs{mec}; \acs{uav}; cloud servers
& 6; 4; 1 \\
Devices; increment; rounds
& 1--200; 1 device/round; 200 \\
Application deadlines: \acs{hrllc}/\acs{mc}/\acs{ic}
& 100/300/500~ms \\
Ground carrier; bandwidth
& 13~GHz; 100~MHz \\
Coverage anchor; bandwidth
& 700~MHz; 10~MHz \\
Aerial access carrier; bandwidth
& Main: 28~GHz; 400~MHz. Ablation: 3.5~GHz; 20~MHz \\
Evaluation seeds; master seed
& 3; 42 \\
State dimension; action slots
& 62; 4 (feasibility-masked) \\
Hidden layers; activation
& 128--128 neurons; ReLU \\
Optimizer; learning rate; loss
& Adam; $10^{-3}$; mean squared error \\
Configured $\gamma$; gradient clipping
& 0.95 (terminal transitions; no bootstrap); 1.0 \\
Replay capacity; minibatch
& 10,000 transitions; 32 \\
Learning; target synchronization
& Every 4 decisions after replay warm-up; every 250 gradient updates \\
Exploration schedule
& $\epsilon$-greedy; nominal $1.0\rightarrow0.05$, linear over 200 training rounds \\
\bottomrule
\end{tabular}
\end{table}




\section{Evaluation}
\label{sec:evaluation}

The proposed framework was evaluated through simulation using 6G-UrbanNetSim The $10\times10$~km deployment combined onboard processing, ground \ac{bs}s with \ac{mec} servers, \ac{uav} computing, and a routed cloud. Active devices increased from 1 to 200 over 200 rounds, alongside growing \ac{ic} payloads. \autoref{tab:experimental-configuration} summarises the experimental configuration. \ac{dqn} family was compared with \ac{ppo}, \ac{a2c}, Q-learning, and Random. Within each campaign, policies shared workload and infrastructure settings; the ablation retrained \ac{ddqn} without \ac{hrllc} latency shaping. Both campaigns used three evaluation seeds; each ablation variant had one training run. 





\subsection{Results and Discussion}
\label{subsec:results-discussion}



%\subsubsection{Training Behaviour}
%\label{sec:eval-training}

% Opening
Consistent \ac{hrllc} offloading requires favourable task outcomes as device load changes. The reward and action histories in~\autoref{fig:train} examine the reward level, its variability, and the execution choices made over 200 evaluation rounds.
%
% Panel (a)
\ac{dqn} and \ac{ddqn} reach high rewards early and sustain the highest levels across most rounds in~\Subfigref{fig:train}{fig:train-reward}, ending at smoothed mean rewards of 3.837 and 3.867, respectively, compared with 3.754 for \ac{ppo} and \ac{a2c}.
%
% Panel (b)
This sustained reward performance is accompanied by relatively low reward variation in~\Subfigref{fig:train}{fig:train-std}. \ac{ddqn} retains particularly consistent terminal outcomes, with a smoothed reward standard deviation of 0.170, compared with 0.437 for \ac{dqn} and 0.514 for \ac{ppo} and \ac{a2c}.
%
% Panel (c)

The action-entropy comparison in~\Subfigref{fig:train}{fig:train-entropy} connects these outcomes to the diversity of offloading decisions. \ac{dqn} and \ac{ddqn} end near 0.649 and 0.640, respectively, reflecting a mixture of execution choices, whereas \ac{ppo} and \ac{a2c} remain at zero because they repeatedly select \ac{mec}. Random exhibits greater action diversity but lower rewards, while \ac{dqn} and \ac{ddqn} combine varied choices with favourable reward outcomes.
%
% Closing
Both proposed variants therefore sustain strong \ac{hrllc} reward performance while using multiple execution options, with \ac{ddqn} providing lower terminal reward variation under the evaluated conditions.



\begin{figure*}[!t]
  \centering
  \subfloat[Mean \acs{hrllc} reward]{\includegraphics[width=\panelwidth]{\paperplotdir/reviewed/Training Convergence - Mean HRLLC Reward.png}\label{fig:train-reward}}
  \hfil
  \subfloat[Reward standard deviation]{\includegraphics[width=\panelwidth]{\paperplotdir/reviewed/Training Stability - Reward Standard Deviation.png}\label{fig:train-std}}
  \hfil
  \subfloat[Action entropy]{\includegraphics[width=\panelwidth]{\paperplotdir/reviewed/Training Exploration - Action Entropy.png}\label{fig:train-entropy}}
  \caption{Reward and action-selection behaviour for \acs{hrllc} during evaluation over 200 rounds.}
  \label{fig:train}
\end{figure*}




\begin{figure*}[!t]
  \centering
  \subfloat[\acs{e2e} delay exceedance]{\includegraphics[width=\quarterwidth]{\paperplotdir/Latency Distribution - E2E Delay Exceedance.png}\label{fig:dist-exceed}}
  \hfil
  \subfloat[Queue delay \acs{cdf}]{\includegraphics[width=\quarterwidth]{\paperplotdir/Latency Distribution - Queue Delay CDF.png}\label{fig:dist-queue}}
  \hfil
  \subfloat[Radio/access delay \acs{cdf}]{\includegraphics[width=\quarterwidth]{\paperplotdir/Latency Distribution - Radio and Access Delay CDF.png}\label{fig:dist-radio}}
  \hfil
  \subfloat[P99/P50 tail inflation]{\includegraphics[width=\quarterwidth]{\paperplotdir/Tail Behaviour - P99 over P50 Inflation.png}\label{fig:tail-infl}}
  \\[2pt]
  \subfloat[Mean \acs{hrllc} latency]{\includegraphics[width=\quarterwidth]{\paperplotdir/Reliability - Mean HRLLC Latency per Round.png}\label{fig:rel-mean-lat}}
  \hfil
  \subfloat[\acs{e2e} P99 per round]{\includegraphics[width=\quarterwidth]{\paperplotdir/Tail Behaviour - End-to-End P99 Latency.png}\label{fig:tail-p99}}
  \hfil
  \subfloat[Queue P99 per round]{\includegraphics[width=\quarterwidth]{\paperplotdir/Delay Composition - Queue Delay P99.png}\label{fig:comp-queue}}
  \hfil
  \subfloat[Handover P99 per round]{\includegraphics[width=\quarterwidth]{\paperplotdir/Delay Composition - Handover Delay P99.png}\label{fig:comp-handover}}
  \\[2pt]
  \subfloat[\acs{hrllc} P99 per round]{\includegraphics[width=\quarterwidth]{\paperplotdir/Reliability - P99 HRLLC Latency per Round.png}\label{fig:rel-p99-lat}}
  \hfil
  \subfloat[P99 versus load]{\includegraphics[width=\quarterwidth]{\paperplotdir/Load Robustness - P99 Latency versus Active Devices.png}\label{fig:load-p99}}
  \hfil
  \subfloat[Deadline success versus load]{\includegraphics[width=\quarterwidth]{\paperplotdir/Load Robustness - Deadline Success versus Active Devices.png}\label{fig:load-success}}
  \hfil
  \subfloat[Radio/access exceedance]{\includegraphics[width=\quarterwidth]{\paperplotdir/Load Robustness - Radio SLA Exceedance versus Active Devices.png}\label{fig:load-radio}}
  \caption{Latency and reliability for \acs{hrllc}}
  \label{fig:latency}
  \label{fig:robustness}
\end{figure*}





\begin{figure*}[!t]
  \centering
  \subfloat[Local/edge execution share]{\edgeshareplot\label{fig:rel-edge}}
  \hfil
  \subfloat[Cloud-offload share]{\includegraphics[width=\quarterwidth]{\paperplotdir/Mobility and Placement - Cloud Offload Share.png}\label{fig:mob-cloud}}
  \hfil
  \subfloat[Mobility violations]{\includegraphics[width=\quarterwidth]{\paperplotdir/Mobility and Placement - Mobility-Induced Violation.png}\label{fig:mob-viol}}
  \hfil
  \subfloat[Local/edge share versus P99]{\includegraphics[width=\quarterwidth]{\paperplotdir/Mobility and Placement - Edge Share versus P99 Delay.png}\label{fig:mob-edge}}
  \\[2pt]
  \subfloat[Cost per byte]{\includegraphics[width=\quarterwidth]{\paperplotdir/Cost Efficiency - Cost per Byte.png}\label{fig:cost-byte}}
  \hfil
  \subfloat[Cost per byte, successful tasks]{\includegraphics[width=\quarterwidth]{\paperplotdir/Cost Efficiency - Cost per Successful Task.png}\label{fig:cost-success}}
  \hfil
  \subfloat[Cost per byte versus P99]{\includegraphics[width=\quarterwidth]{\paperplotdir/Cost Efficiency - Cost versus P99 Latency.png}\label{fig:cost-p99}}
  \hfil
  \subfloat[Useful bits/J versus P99]{\includegraphics[width=\quarterwidth]{\paperplotdir/Tail Behaviour - P99 versus Energy Efficiency.png}\label{fig:tail-eff}}
  \caption{Offloading and efficiency of \acs{hrllc}}
  \label{fig:efficiency}
  \label{fig:reliability}
\end{figure*}



\begin{figure*}[!t]
  \centering
  \subfloat[End-to-end delay exceedance]{\includegraphics[width=\panelwidth]{\paperplotdir/Ablation - End-to-End Delay Exceedance.png}\label{fig:ablation-exceedance}}
  \hfil
  \subfloat[P99 latency versus active devices]{\includegraphics[width=\panelwidth]{\paperplotdir/Ablation - P99 Latency versus Active Devices.png}\label{fig:ablation-load-p99}}
  \hfil
  \subfloat[Mean delay components]{\includegraphics[width=\panelwidth]{\paperplotdir/Ablation - Delay Component Breakdown.png}\label{fig:ablation-components}}
  \\[2pt]
  \subfloat[Execution-tier share versus active devices]{\includegraphics[width=\panelwidth]{\paperplotdir/Ablation - Execution Tier Share versus Active Devices.png}\label{fig:ablation-placement}}
  \hfil
  \subfloat[Radio and access delay exceedance]{\includegraphics[width=\panelwidth]{\paperplotdir/Ablation - Radio SLA Exceedance versus Active Devices.png}\label{fig:ablation-radio}}
  \hfil
  \subfloat[P99 latency versus energy efficiency]{\includegraphics[width=\panelwidth]{\paperplotdir/Ablation - P99 Latency versus Energy Efficiency.png}\label{fig:ablation-energy}}
  \caption{Ablation study of \acs{ddqn} with and without latency shaping.}
  \label{fig:ablation}
\end{figure*}




\subsubsection{Reliability}
\label{sec:eval-latency}

% Opening
Latency-sensitive urban tasks require fast completion and dependable service as device demand increases. The latency distributions, delay components, and load-dependent results in~\autoref{fig:latency} evaluate these requirements for \ac{hrllc} offloading.
%
% Panel (a)
The \ac{e2e} distribution in~\Subfigref{fig:latency}{fig:dist-exceed} gives pooled P99 values of 1.053\,ms for \ac{dqn} and 1.062\,ms for \ac{ddqn}, compared with 2.037\,ms for \ac{ppo} and \ac{a2c}, representing approximately 48\% lower tail latency under both proposed variants.
%
% Panel (b)
This favourable latency performance is accompanied by limited queueing in~\Subfigref{fig:latency}{fig:dist-queue}: 91.72\% of \ac{dqn} tasks and 89.47\% of \ac{ddqn} tasks experience zero queue delay, compared with 82.96\% under \ac{ppo} and \ac{a2c}.
%
% Panel (c)
The radio/access distributions in~\Subfigref{fig:latency}{fig:dist-radio} provide a similar comparison across the complete \ac{hrllc} workload, including local execution, with P99 values of 0.0197\,ms for \ac{dqn} and 0.0350\,ms for \ac{ddqn}, compared with 1.021\,ms for \ac{ppo} and \ac{a2c}.
%
% Panel (d)

The round-level tail-inflation history in~\Subfigref{fig:latency}{fig:tail-infl} further identifies \ac{ddqn}'s consistent behaviour, with its smoothed P99/P50 ratio peaking at 1.87.
%
% Panel (e)
Mean latency in~\Subfigref{fig:latency}{fig:rel-mean-lat} distinguishes the two proposed variants: \ac{ddqn} records a campaign mean of 1.025\,ms, compared with 1.620\,ms for \ac{dqn}, whose higher mean reflects isolated slow completions.
%
% Panel (f)
The round-level \ac{e2e} P99 in~\Subfigref{fig:latency}{fig:tail-p99} reinforces \ac{ddqn}'s consistency, reaching a maximum smoothed seed-mean value of 1.899\,ms, compared with approximately 121\,ms for \ac{ppo} and \ac{a2c}.
%
% Panel (g)
Queue P99 in~\Subfigref{fig:latency}{fig:comp-queue} links temporary \ac{e2e} increases to queue build-up,
%
% Panel (h)
while the recorded access/handover waits in~\Subfigref{fig:latency}{fig:comp-handover} identify another contribution to isolated delays.
%
% Panel (i)
The \ac{hrllc} P99 history in~\Subfigref{fig:latency}{fig:rel-p99-lat} traces the same round-level tail pattern.
%
% Panel (j)

Under increasing load, \Subfigref{fig:latency}{fig:load-p99} gives a mean per-seed P99 of 1.103\,ms for \ac{ddqn} at the last plotted load of 197 devices, compared with 32.342\,ms for \ac{dqn}, identifying \ac{ddqn}'s additional high-load tail-latency advantage.
%
% Panel (k)
Application success in~\Subfigref{fig:latency}{fig:load-success} connects these delays to service delivery: \ac{ddqn} meets the 100\,ms deadline in all 17,105 attempts, compared with 17,101 under \ac{dqn} and 17,100 under \ac{ppo} and \ac{a2c}.
%
% Panel (l)
The separate radio/access comparison in~\Subfigref{fig:latency}{fig:load-radio} also favours the proposed variants near the end of the ramp, with exceedance rates of 0.463\% for \ac{ddqn} and 0.612\% for \ac{dqn} at the 0.1\,ms threshold, compared with approximately 2.997\% for \ac{ppo} and \ac{a2c}.
%
% Closing
Both proposed variants therefore combine low pooled tail latency with high observed application-deadline success, while \ac{ddqn} provides more consistent tail behaviour under the evaluated load conditions.





\subsubsection{Offloading and Efficiency}
\label{subsubsec:results-cost-energy}

% Opening
Timely \ac{hrllc} offloading depends on the execution tiers selected and the resources consumed by those decisions. The offloading, mobility, cost, and energy comparisons in~\autoref{fig:efficiency} relate these execution choices to the latency performance of the proposed variants.
%
% Panel (a)
\ac{dqn} and \ac{ddqn} maintain 100\% local/edge execution share throughout the evaluation in~\Subfigref{fig:efficiency}{fig:rel-edge}, matching \ac{ppo} and \ac{a2c}. This metric includes local execution. Both proposed variants combine local and offloaded execution, with campaign-wide local shares of 37.40\% for \ac{dqn} and 36.11\% for \ac{ddqn}, whereas \ac{ppo} and \ac{a2c} allocate their \ac{hrllc} tasks exclusively to \ac{mec}.
%
% Panel (b)
This offloading pattern is reflected in the zero cloud-offload share of both proposed variants in~\Subfigref{fig:efficiency}{fig:mob-cloud}. Across the complete campaign, Q-learning and Random allocate 14.84\% and 33.65\% of their \ac{hrllc} tasks to the cloud, respectively.
%
% Panel (c)
The mobility comparison in~\Subfigref{fig:efficiency}{fig:mob-viol} records campaign-wide violation rates of 16.71\% for \ac{dqn} and 17.80\% for \ac{ddqn}, compared with 36.61\% for \ac{ppo} and \ac{a2c}. These lower recorded rates accompany the proposed variants' mixture of local and offloaded execution. The dwell-time check applies to \ac{mec} and aerial offloads, so the aggregate rates also reflect the allocation mix.
%
% Panel (d)

The relationship between local/edge execution share and latency in~\Subfigref{fig:efficiency}{fig:mob-edge} connects this allocation behaviour to task completion time. At 100\% mean per-seed local/edge execution share, \ac{dqn} and \ac{ddqn} record mean per-seed P99 values of 1.052 and 1.058\,ms, respectively, compared with 2.037\,ms for \ac{ppo} and \ac{a2c} at the same local/edge execution share.
%
% Panel (e)
This low tail latency is accompanied by low recorded unit cost in~\Subfigref{fig:efficiency}{fig:cost-byte}. At the smoothed terminal point, \ac{dqn} and \ac{ddqn} record approximately $2.754\times10^{-10}$ and $2.805\times10^{-10}$\,USD/byte, respectively, compared with $3.049\times10^{-10}$\,USD/byte for Q-learning.
%
% Panel (f)
Restricting the calculation to timely tasks preserves this cost advantage in~\Subfigref{fig:efficiency}{fig:cost-success}. The corresponding terminal values are $2.748\times10^{-10}$\,USD/byte for \ac{dqn} and $2.805\times10^{-10}$\,USD/byte for \ac{ddqn}, compared with $2.955\times10^{-10}$\,USD/byte for Q-learning. The unchanged \ac{ddqn} value accompanies its complete observed application-deadline success.
%
% Panel (g)
The campaign cost--latency comparison in~\Subfigref{fig:efficiency}{fig:cost-p99} places both proposed variants at the low-cost, low-latency end of the evaluated operating points. Mean per-seed unit cost is $2.759\times10^{-10}$\,USD/byte for \ac{dqn} and $2.822\times10^{-10}$\,USD/byte for \ac{ddqn}, compared with $3.545\times10^{-10}$\,USD/byte for Q-learning, while both proposed variants retain mean per-seed P99 values near 1.05\,ms.
%
% Panel (h)

The energy comparison in~\Subfigref{fig:efficiency}{fig:tail-eff} identifies the trade-off accompanying this latency performance. Mean task efficiency, averaged across seeds, is $2.901\times10^{6}$~useful bits/J for \ac{dqn} and $2.900\times10^{6}$~useful bits/J for \ac{ddqn}, compared with approximately $4.176\times10^{6}$~useful bits/J for \ac{ppo} and \ac{a2c}. Their higher efficiency accompanies a mean per-seed P99 of approximately 2.037\,ms, compared with approximately 1.05\,ms under the proposed variants.
%
% Closing
Both proposed variants therefore combine exclusive local/edge execution, lower recorded mobility-violation rates than \ac{ppo} and \ac{a2c}, and low unit cost with short tail latency. The energy comparison defines the resource trade-off accompanying this latency-focused behaviour under the evaluated urban \ac{hrllc} workload.






\subsubsection{Ablation Study}
\label{sec:eval-ablation}

% Opening
The \ac{hrllc} latency contribution encourages \ac{ddqn} to favour shorter task-completion times. The matched ablation in~\autoref{fig:ablation} examines this contribution by comparing full \ac{ddqn} with a retrained variant whose latency shaping is disabled, while retaining reward normalisation and the deadline-miss penalty. The comparison uses one matched training replicate and three evaluation seeds.
%
% Panel (a)
Full \ac{ddqn} maintains a shorter E2E exceedance tail in~\Subfigref{fig:ablation}{fig:ablation-exceedance}. Its mean per-seed completed-task P99 is 1.054\,ms, compared with 2.297\,ms when latency shaping is disabled, representing approximately 54.1\% lower tail latency.
%
% Panel (b)
This latency advantage persists across increasing device demand in~\Subfigref{fig:ablation}{fig:ablation-load-p99}, with full \ac{ddqn} recording lower mean P99 in every 25-device load bin. In the 176--200-device bin, full and ablated \ac{ddqn} record 1.073 and 2.035\,ms, respectively, showing that the benefit extends to the highest evaluated load range.
%
% Panel (c)

The component comparison in~\Subfigref{fig:ablation}{fig:ablation-components} links this advantage to shorter communication and waiting delays. Full \ac{ddqn} records mean transmission, propagation, and queueing delays of 0.0076, 0.0089, and 0.0043\,ms, compared with 0.0111, 0.0718, and 0.0296\,ms under the ablated variant, while processing remains near 1.000\,ms in both cases. Mean handover delay is 0.0053\,ms under full \ac{ddqn} and 0.3193\,ms under the ablated variant. The latter mean is strongly influenced by one 5,000\,ms event, which accounts for 91.56\% of its handover-delay sum.
%
% Panel (d)
Execution-tier selection changes alongside these delay components in~\Subfigref{fig:ablation}{fig:ablation-placement}. Full \ac{ddqn} retains the \ac{hrllc} workload within local and \ac{mec} resources, with zero cloud allocation, whereas disabling latency shaping produces a mean per-seed cloud share of 22.42\%. The ablated allocation varies substantially with load; its cloud share falls below 0.4\% in the final bin while its P99 remains higher, showing that the latency difference is not restricted to periods of substantial cloud use.
%
% Panel (e)

The radio/access comparison in~\Subfigref{fig:ablation}{fig:ablation-radio} provides further support for the latency contribution. In the final load bin, the share of \ac{hrllc} tasks exceeding the 0.1\,ms radio/access threshold is 0.636\% under full \ac{ddqn}, compared with 2.593\% when latency shaping is disabled.
%
% Panel (f)
The energy--latency operating points in~\Subfigref{fig:ablation}{fig:ablation-energy} identify the resource trade-off accompanying this improvement. Mean per-seed task efficiency is 3.340\,Mbit/J under full \ac{ddqn} and 3.846\,Mbit/J under the ablated variant, whose greater efficiency accompanies the longer E2E tail. Full \ac{ddqn} also meets the 100\,ms application deadline in all 17,105 attempts, compared with 17,104 under the ablated variant.
%
% Closing
The matched comparison therefore supports the role of latency reward shaping in obtaining shorter \ac{hrllc} completion tails across the evaluated load range. Full \ac{ddqn} combines lower tail latency and radio/access exceedance with complete observed application-deadline success, while the energy comparison makes the associated efficiency trade-off explicit.



\section{Conclusion}
\label{sec:conclusion}

This paper developed a reliability-aware task-offloading framework for heterogeneous urban workloads across the 6G computing continuum. The framework uses service-differentiated rewards to prioritise \ac{hrllc} latency and deadline requirements when selecting execution tiers. We developed \textit{6G-UrbanNetSim} simulator to evaluate offloading policies across onboard, ground-edge, aerial, and cloud resources under increasing device demand.
%
The evaluation showed that the proposed approach supported low \ac{hrllc} tail latency and high application-deadline success. The comparative evaluation and reward-shaping ablation further supported the contribution of latency-aware task offloading, while identifying the associated resource trade-offs.
%
The evaluation covers one simulated urban deployment in which device count and \acs{ic} payload size increase together. The reward ablation uses one matched training pair evaluated on three seeds, so it measures the effect for that pair. Independent training replicates are needed to estimate how much the effect varies across runs.





\bibliographystyle{ieeetr}
\bibliography{references}

\clearpage
\onecolumn
\appendices
\renewcommand{\theequation}{\Alph{section}.\arabic{equation}}
\renewcommand{\theHequation}{\Alph{section}.\arabic{equation}}
\setcounter{section}{0}
\setcounter{equation}{0}
\setcounter{table}{0}
\renewcommand*{\theHtable}{appendix.\arabic{table}}

\section{State Encoding and Fixed Normalization}
\label{sec:appendix-state-encoding}

The offloading agent receives the 62-entry observation defined in~\eqref{eq:observation-encoding}.
The first 22 entries describe the task, device, and network conditions, as shown in Table~\ref{tab:state-context-features}.
The remaining 40 entries are arranged in four blocks of ten features, one block for each action slot.
Table~\ref{tab:state-action-features} gives the order of features within each block.
For slot \(a=0,\ldots,3\), feature \(k=1,\ldots,10\) show position \(22+10a+k\) in the observation.
Slot \(a=0\) corresponds to onboard execution.
%
To normalize a value \(x\), the encoder divides it by a positive scale \(b\) and clips the result to the interval \([0,1]\).

\begin{equation}
\mathcal N(x;b)=\operatorname{clip}(x/b,0,1).
\label{eq:appendix-fixed-normalization}
\end{equation}
The encoder uses fixed scales throughout training and evaluation.

The first four context features record the task payload in kilobytes and the values of the legacy \acs{cpu}, memory, and storage demand fields.
The encoder divides all four values by \(M_D=2.5D_{\mathrm{cfg}}\), where \(D_{\mathrm{cfg}}\) is the fixed value of the \texttt{data\_per\_device\_kb} environment.
The payload in the first feature is the size of the individual task.
The scaling applies only to the observation inputs.
Compute-rate demand, available capacity, and battery energy retain their physical units in the admission, execution, and reward calculations.

The encoder scales the handover count by the configured maximum number of devices, \(N_{\max}\).
For the distance features, it uses \(L=\max\{1,\sqrt{X_{\max}^2+Y_{\max}^2}\}\), where \(X_{\max}\) and \(Y_{\max}\) are the upper bounds of the city coordinates in metres.
In the evaluated configuration, \(D_{\mathrm{cfg}}=400\), \(M_D=1000\), \(N_{\max}=200\), and \(X_{\max}=Y_{\max}=10{,}000\)~m, giving \(L\approx14{,}142\)~m.
For feature 17, the mobility observer computes the mean distance of the active devices from their position centroid and rounds it to two decimal places before normalization.

\begin{table}[!htbp]
\centering
\caption{Ordered context features. Indices refer to the first 22 observation positions.}
\label{tab:state-context-features}
\small
\setlength{\tabcolsep}{5pt}
\renewcommand{\arraystretch}{1.12}
\begin{tabular}{@{}r p{0.51\linewidth} p{0.36\linewidth}@{}}
\toprule
Index & Input quantity & Encoded value \\
\midrule
1 & Task payload, numerical kilobyte value & \(\mathcal N(\text{payload};M_D)\) \\
2 & Legacy \acs{cpu} demand, \texttt{cpu\_demand} & \(\mathcal N(\text{legacy CPU demand};M_D)\) \\
3 & Legacy memory demand, \texttt{memory\_demand} & \(\mathcal N(\text{legacy memory demand};M_D)\) \\
4 & Legacy storage demand, \texttt{storage\_demand} & \(\mathcal N(\text{legacy storage demand};M_D)\) \\
5 & Service priority (\(p_i\)): 1 for \acs{hrllc}, 2 for \acs{mc}, 3 for \acs{ic} & \(\mathcal N(p_i;3)\) \\
6 & Application deadline, in milliseconds & \(\mathcal N(\text{deadline};1000)\) \\
7 & Device-to-serving-\acs{bs} distance (\(D_i^{\mathrm{BS}}\)), in metres & \(\mathcal N(D_i^{\mathrm{BS}};L)\) \\
8 & Received signal level (\(S_i^{\mathrm{rx}}\)), in dBm & \(\mathcal N(|S_i^{\mathrm{rx}}|;120)\) \\
9 & Device has onboard compute capability & Binary indicator \\
10 & Device handover delay, in milliseconds & \(\mathcal N(\text{handover delay};500)\) \\
11 & Device movement speed, in metres per second & \(\mathcal N(\text{speed};30)\) \\
12 & Predicted dwell time, in milliseconds & \(\mathcal N(\text{dwell time};60{,}000)\) \\
13 & Device handover-risk value & \(\mathcal N(\text{handover risk};1)\) \\
14 & Device remaining-battery fraction \(\beta_i\) & Clipped to \([0,1]\) \\
15 & Number of handovers in the current mobility step & \(\mathcal N(\text{handover count};N_{\max})\) \\
16 & Sum of handover delays in the current mobility step, in milliseconds & \(\mathcal N(\text{handover-delay sum};1000)\) \\
17 & Mean distance of active devices from their position centroid, in metres & \(\mathcal N(\text{position dispersion};L)\) \\
18 & Device horizontal coordinate (\(x_i\)), in metres & \(\mathcal N(x_i;X_{\max})\) \\
19 & Device battery below the configured low-battery fraction & \(\mathbf 1\{\beta_i<0.30\}\) \\
20 & Service indicator: \acs{hrllc} & \(\mathbf 1\{p_i=1\}\) \\
21 & Service indicator: \acs{mc} & \(\mathbf 1\{p_i=2\}\) \\
22 & Service indicator: \acs{ic} & \(\mathbf 1\{p_i=3\}\) \\
\bottomrule
\end{tabular}
\end{table}

Each action block describes the conditions the task would face under that execution option.
Current utilization \(u_{ia}\) measures the fraction of capacity in use before the task arrives.
Projected utilization \(\widehat u_{ia}\) estimates the fraction that would be in use after adding the task.
For onboard execution, both values are based on \acs{cpu} utilization.
For a remote option, the encoder uses the highest utilization among \acs{cpu}, memory, and storage at the selected server.
Current and projected headroom give the corresponding fractions of capacity that are free.

For a remote option, \(g_{ia}\) is the congestion percentage reported for the selected server.
The encoder divides this percentage by 100 and clips the result to \([0,1]\).

Compute adequacy \(\kappa_{ia}\) compares the execution option's total compute capacity \(C_{ia,\mathrm{cpu}}\) with the task's required compute rate \(\rho_i\).
Both are measured in cycles per second.
When both values are positive, the encoder divides capacity by demand and clips the result to \([0,1]\).
The score is zero when capacity is zero or negative.
It is one when capacity is positive and demand is zero or negative.

\begin{table}[H]
\centering
\caption{Ordered features repeated for each of the four action slots.}
\label{tab:state-action-features}
\small
\setlength{\tabcolsep}{5pt}
\renewcommand{\arraystretch}{1.12}
\begin{tabular}{@{}r p{0.51\linewidth} p{0.36\linewidth}@{}}
\toprule
Block index & Input quantity & Encoded value \\
\midrule
1 & Physical admission and route-feasibility mask \(m_{ia}\) & Binary indicator \\
2 & Projected utilization & \(\operatorname{clip}(\widehat u_{ia},0,1)\) \\
3 & Resolved-server congestion percentage & \(\mathcal N(g_{ia};100)\) \\
4 & Projected free headroom & \(1-\operatorname{clip}(\widehat u_{ia},0,1)\) \\
5 & Current free headroom & \(\operatorname{clip}(1-u_{ia},0,1)\) \\
6 & Task-relative compute adequacy & \(\kappa_{ia}\) \\
7 & Remaining-battery fraction of the resolved server (\(\beta_{ia}^{S}\)) & Clipped to \([0,1]\) \\
8 & Resolved-server battery below the low-battery fraction & \(\mathbf 1\{\beta_{ia}^{S}<0.30\}\) \\
9 & Action availability & Binary indicator, defined below \\
10 & Action is onboard execution & \(\mathbf 1\{a=0\}\) \\
\bottomrule
\end{tabular}
\end{table}

For onboard execution, the availability flag is one when local computing is enabled and the device supports it.
The feasibility mask indicates whether the device can accept the particular task under the applicable resource and battery checks.
For a remote option, the availability flag is one when the server-selection procedure finds an operational, reachable server that can accept the task.

To compute the server's remaining-battery fraction, the encoder divides the remaining energy by the battery capacity and clips the result to \([0,1]\).
If there is no selected server or its battery capacity is zero or negative, the encoder sets the fraction to one.
Onboard execution uses this default and records zero server congestion.
The device's own battery fraction appears in feature 14 of Table~\ref{tab:state-context-features}.


\end{document}
